\chapter*{Abstract}

\section*{\tituloPortadaEngVal}

Progress in an individual sport like climbing is often slow and marked by long plateaus. This poses a significant challenge in maintaining a climber's motivation and consistency. Adding to this problem is the fact that tracking progress and achievements is done manually in most cases, relying on memory or physical notebooks. To solve this issue, this project is built upon two fundamental pillars: the context of climbing gyms and the use of gamification techniques applied to sports.

The value proposition consists of a mobile application that digitizes the sports management of climbing gyms and provides climbers with the ability to track their activity. To tackle the problem of demotivation, the tool incorporates gamification mechanics based on the PB system (Points, Badges, and Leaderboards). These elements transform the training routine into a more dynamic experience, allowing users to visualize their progress and encouraging a healthy competition that reinforces a sense of belonging to the climbing community.

The project has been developed under a UCD (User-Centered Design) approach, ensuring that the requirement definition, improvement proposals, and testing stem directly from the actual needs of the climbers. This guarantees that the built solution provides real value and adapts to the needs of the end users.

\section*{Keywords}

\noindent Indoor climbing, gamification, climbing gyms, motivation, PBL system, user-centered design, mobile application, sports tracking.



